\chapter{Introduzione}

La donazione di sangue è un atto volontario e altruistico che costituisce uno
dei pilastri del sistema sanitario moderno. Il sangue, come piastrine e plasma,
non sono sostanze sintetizzabili artificialmente in laboratorio: 
la loro disponibilità dipende in modo esclusivo dalla scelta dei cittadini di donare.
In Italia il sistema trasfusionale si regge su una rete capillare di associazioni di 
volontariato e su centri ospedalieri regionali che gestiscono raccolta, lavorazione e
distribuzione degli emocomponenti.

Il Friuli Venezia Giulia presenta, nel suo insieme, caratteristiche degne di
nota per le analisi che seguiranno: una popolazione mediamente anziana (48,8 anni,
1,7 in più della media nazionale~\cite{etàmediaFVG}) e una forte presenza di
volontariato. Il presente lavoro, tuttavia, non copre l'intera regione. In Friuli Venezia
Giulia operano tre Dipartimenti Interaziendali di Medicina Trasfusionale,
ciascuno dei quali fa riferimento ad un territorio specifico, gestendo
le operazioni di donazione e raccolta tra presidi ospedalieri principali e 
autoemoteche che operano nei comuni minori:
quello dell'Azienda Sanitaria Universitaria Giuliano Isontina (ASUGI), 
competente su Trieste e Gorizia; quello
dell'Azienda Sanitaria Universitaria Friuli Centrale (ASUFC), su Udine; e
quello dell'Azienda Sanitaria Friuli Occidentale (ASFO), su Pordenone. I dati
disponibili per questo lavoro sono quelli del \textbf{DIMT dell'ASUGI}.

Le province di Trieste e Gorizia rappresentano circa un quinto del bacino
donazionale regionale e non ne riproducono le caratteristiche: la realtà
sociale ed economica dell'area giuliano-isontina, e in particolare del
capoluogo triestino, si differenzia per molti aspetti --- demografici,
economici e organizzativi --- da quella di Udine e Pordenone. Le conclusioni
di questo lavoro vanno perciò lette come relative al territorio
effettivamente osservato, e non estese senza verifica al resto della regione.

Comprendere i fattori che influenzano la frequenza e la continuità delle
donazioni è rilevante sia sul piano scientifico sia su quello operativo:
permette di identificare i profili a rischio di abbandono, di pianificare
campagne di reclutamento mirate e di anticipare fluttuazioni nell'offerta di
emocomponenti. Il presente lavoro approfondisce questi aspetti analizzando un
dataset longitudinale di \textbf{diciassette anni} di attività donazionale
nelle province di Trieste e Gorizia, dal 2009 al 2025, con strumenti di
statistica descrittiva e di modellazione regressiva.

Una precisazione di metodo, perché orienta la lettura di tutto ciò che segue.
Il calo delle donazioni è un fatto noto e ripetutamente osservato; è invece
raramente scomposto. Un sistema che raccoglie meno sangue in un territorio che
perde abitanti può essere in erosione oppure semplicemente più piccolo, e le
due diagnosi portano a decisioni opposte. Il \cref{ch:eda} affronta la
distinzione esplicitamente, rapportando le donazioni alla popolazione
residente, e il risultato che ne esce condiziona la lettura dei capitoli
successivi.

Il contributo della tesi è duplice. Sul piano applicativo, quantifica per le
province di Trieste e Gorizia un fenomeno spesso descritto solo
qualitativamente e ne individua i punti
di intervento --- la prima ripetizione dei donatori giovani e il rientro della
fascia di mezza età --- laddove il margine sui donatori già attivi risulta
invece vincolato in alto da regole cliniche. Sul piano metodologico, mostra che
tre scelte apparentemente tecniche cambiano le conclusioni: quale colonna
identifica il donatore, se un campo amministrativo incompleto sia un requisito
o una covariata, e su quale unità di analisi si misuri un conteggio. Ogni
figura e ogni numero riportati sono tracciabili al comando che li ha prodotti.